
\section{Problem Formulation and A2UI Primer}
\label{sec:problem}

We study A2UI-based Generative UI: given a system instruction, a dialogue history, and the current user message, the model must produce a unified assistant response that contains natural language and, when appropriate, a structured A2UI message sequence. Unlike approaches that ask the model to generate HTML, JavaScript, or framework-specific code, A2UI is a declarative protocol in which the model emits structured messages, and the client renders them using a trusted component catalog. This separation is important for our setting: it makes UI generation safer, more portable across rendering environments, and easier to validate automatically. In this paper, we instantiate all data construction, rendering, and evaluation against A2UI v0.8, the current stable public version.

At a high level, A2UI v0.8 organizes interaction through four message types. \texttt{surfaceUpdate} defines or updates UI components, \texttt{dataModelUpdate} updates application state, \texttt{beginRendering} signals the client to render a surface, and \texttt{deleteSurface} removes an existing surface. 

There are several challenges in A2UI generation. The first one is \textbf{protocol validity}: an output may be syntactically well-formed JSON yet still violate message, reference, typing, or renderability constraints. The second is \textbf{interaction construction}: even a protocol-valid UI may still use the wrong widget type, fail to ground visible choices in the assistant text, or mishandle state updates across turns. The third challenge is \textbf{user-facing quality}: an output may be structurally correct and functionally plausible, yet still add little value beyond plain text, feel abrupt in context, or impose unnecessary cognitive load on the user.

This decomposition directly motivates the rest of the paper. In Section~\ref{sec:data}, we construct an A2UI-grounded corpus that teaches the model how to generate protocol-compliant UI. In Section~\ref{sec:benchmark}, we design a benchmark that evaluates model's ability when facing these challenges. In Section~\ref{sec:exp}, we show that learning under this formulation benefits from a two-stage training pipeline: supervised fine-tuning first stabilizes the response format and basic text--UI grounding, while reinforcement learning further improves the quality of executable interaction.